\documentclass[ctcc_full_documentation.tex]{subfiles}
\begin{document}
\chapter{Delay Costs Calculation Documentation}

\section{Introduction}

The \texttt{delay\_costs.py} script calculates the costs incurred during the delay period (permitting/regulatory delays) for a transmission line project. These costs include legal, administrative, labor, materials, regulatory, public relations, project management, and miscellaneous expenses.

\section{Method Overview}

\subsection{Purpose}

This script calculates delay costs by:
\begin{enumerate}
    \item Loading annual delay cost components from YAML configuration
    \item Summing all annual delay cost components
    \item Multiplying by the number of delay years to get total nominal cost
    \item Calculating present value of annual delay costs
\end{enumerate}

\subsection{Calculation Approach}

The delay cost calculation follows this structure:

\begin{align}
\text{Total Yearly Delay Cost} &= \sum_{\text{components}} \text{Annual Component Cost} \\
\text{Total Delay Cost} &= \text{Total Yearly Delay Cost} \times \text{Delay Years} \\
\text{Present Value} &= \sum_{t=1}^{D} \frac{\text{Total Yearly Delay Cost}}{(1 + r)^t}
\end{align}

where $D$ is the number of delay years and $r$ is the real WACC discount rate.

\subsection{Inputs and Outputs}

\textbf{Inputs:}
\begin{itemize}
    \item Annual delay cost components (legal, admin, labor, materials, regulatory, public relations, project management, miscellaneous)
    \item Delay years (from project timeline)
    \item Financing parameters (WACC for present value)
\end{itemize}

\textbf{Outputs:}
\begin{itemize}
    \item Total yearly delay cost
    \item Total nominal delay cost
    \item Present value of delay costs
\end{itemize}

\subsection{Integration with Other Scripts}

This script integrates with:
\begin{itemize}
    \item \texttt{smart\_loaders.py} --- \texttt{load\_delay\_costs}, \texttt{load\_financing\_details}, \texttt{load\_project\_technical\_details} (centralized; local wrapper \texttt{load\_project\_technical\_details()} returns \texttt{delay\_years} only)
    \item \texttt{financial\_utils.py} --- \texttt{calculate\_present\_value}
    \item \texttt{smart\_output.py} --- \texttt{CTCCOutputManager} (\texttt{add\_delay\_costs}, \texttt{write\_batch\_summary})
\end{itemize}

\section{Key Functions}

\subsection{\texttt{load\_project\_technical\_details()}}

Local wrapper that loads project technical details via \texttt{smart\_loaders.load\_project\_technical\_details\_centralized} and returns only \texttt{delay\_years}.

\textbf{Returns:}
\begin{itemize}
    \item \texttt{delay\_year} (float): Number of delay years
\end{itemize}

\subsection{\texttt{main()}}

Main execution function that orchestrates the delay cost calculation and output.

\textbf{Workflow:}
\begin{enumerate}
    \item Loads delay cost components from YAML
    \item Loads delay years from project technical details
    \item Sums all annual delay cost components
    \item Calculates total nominal cost (annual cost × delay years)
    \item Calculates present value of annual delay costs
    \item Displays results and writes to CSV
\end{enumerate}

\textbf{Key Logic:}
\begin{lstlisting}
# Load annual delay cost components
legal = delay_cost_df["annual_delay_costs"]["legal"]
admin = delay_cost_df["annual_delay_costs"]["admin"]
labor = delay_cost_df["annual_delay_costs"]["labor"]
material_and_equipment = delay_cost_df["annual_delay_costs"]["material_and_equipment"]
regulatory = delay_cost_df["annual_delay_costs"]["regulatory"]
public_relations = delay_cost_df["annual_delay_costs"]["public_relations"]
project_management = delay_cost_df["annual_delay_costs"]["project_management"]
miscellaneous = delay_cost_df["annual_delay_costs"]["miscellaneous"]

# Sum components
total_yearly_delay_cost = (legal + admin + labor + material_and_equipment 
                          + regulatory + public_relations + project_management 
                          + miscellaneous)

# Calculate total nominal cost
total_delay_cost = total_yearly_delay_cost * delay_year

# Calculate present value (payments from year 1 for delay_year years)
total_delay_cost_pv = calculate_present_value(
    total_yearly_delay_cost,
    financing.wacc_real,
    delay_year,
)
\end{lstlisting}

\section{Example}

The following example demonstrates a typical usage of the delay cost calculation:

\begin{lstlisting}
# Example: 2-year delay period
# 
# Inputs (from YAML):
#   - Legal: $50,000/year
#   - Admin: $30,000/year
#   - Labor: $100,000/year
#   - Materials: $20,000/year
#   - Regulatory: $40,000/year
#   - Public Relations: $25,000/year
#   - Project Management: $80,000/year
#   - Miscellaneous: $15,000/year
#   - Delay years: 2
#   - Real WACC: 3%
#
# Calculation:
#   Total yearly delay cost = $50k + $30k + $100k + $20k + $40k 
#                           + $25k + $80k + $15k = $360,000/year
#   Total nominal cost = $360,000 * 2 = $720,000
#   Present value = $360,000 * PV annuity (2 years, 3%) = $697,000
#
# Output:
#   - Total yearly delay cost: $360,000
#   - Total nominal cost: $720,000
#   - Present value: $697,000
\end{lstlisting}

\textbf{Integration Example:}

\begin{lstlisting}
# When called from ctcc.py:
csv_manager = CTCCOutputManager()
results = {
    "total_nominal": total_delay_cost,
    "total_afudc": 0,  # Delay costs are not AFUDC-eligible
    "total_pv": total_delay_cost_pv,
}
csv_manager.add_delay_costs(results)
csv_manager.write_batch_summary()
\end{lstlisting}

\section{Key Implementation Details}

\subsection{Delay Cost Components}

The script includes eight delay cost components:
\begin{itemize}
    \item \textbf{Legal}: Legal fees and expenses
    \item \textbf{Admin}: Administrative costs
    \item \textbf{Labor}: Labor costs during delay
    \item \textbf{Material and Equipment}: Material and equipment holding costs
    \item \textbf{Regulatory}: Regulatory compliance costs
    \item \textbf{Public Relations}: Public relations and community engagement
    \item \textbf{Project Management}: Project management overhead
    \item \textbf{Miscellaneous}: Other delay-related expenses
\end{itemize}

\subsection{AFUDC Treatment}

Delay costs are not AFUDC-eligible because they are operating expenses during the delay period, not construction costs. AFUDC applies only to construction work in progress (CWIP).

\subsection{Present Value Timing}

Delay costs are incurred annually during the delay period (years 1 to delay\_year). The present value calculation discounts these annual payments back to the base year.

\subsection{Simplified Structure}

This script has a simpler structure than other cost scripts because delay costs are straightforward: annual costs multiplied by delay years, with no complex terrain adjustments or component-level calculations.

\end{document}

